% Options for packages loaded elsewhere
\PassOptionsToPackage{unicode}{hyperref}
\PassOptionsToPackage{hyphens}{url}
%
\documentclass[
]{article}
\usepackage{amsmath,amssymb}
\usepackage{iftex}
\ifPDFTeX
  \usepackage[T1]{fontenc}
  \usepackage[utf8]{inputenc}
  \usepackage{textcomp} % provide euro and other symbols
\else % if luatex or xetex
  \usepackage{unicode-math} % this also loads fontspec
  \defaultfontfeatures{Scale=MatchLowercase}
  \defaultfontfeatures[\rmfamily]{Ligatures=TeX,Scale=1}
\fi
\usepackage{lmodern}
\ifPDFTeX\else
  % xetex/luatex font selection
\fi
% Use upquote if available, for straight quotes in verbatim environments
\IfFileExists{upquote.sty}{\usepackage{upquote}}{}
\IfFileExists{microtype.sty}{% use microtype if available
  \usepackage[]{microtype}
  \UseMicrotypeSet[protrusion]{basicmath} % disable protrusion for tt fonts
}{}
\makeatletter
\@ifundefined{KOMAClassName}{% if non-KOMA class
  \IfFileExists{parskip.sty}{%
    \usepackage{parskip}
  }{% else
    \setlength{\parindent}{0pt}
    \setlength{\parskip}{6pt plus 2pt minus 1pt}}
}{% if KOMA class
  \KOMAoptions{parskip=half}}
\makeatother
\usepackage{xcolor}
\usepackage[margin=1in]{geometry}
\usepackage{color}
\usepackage{fancyvrb}
\newcommand{\VerbBar}{|}
\newcommand{\VERB}{\Verb[commandchars=\\\{\}]}
\DefineVerbatimEnvironment{Highlighting}{Verbatim}{commandchars=\\\{\}}
% Add ',fontsize=\small' for more characters per line
\usepackage{framed}
\definecolor{shadecolor}{RGB}{248,248,248}
\newenvironment{Shaded}{\begin{snugshade}}{\end{snugshade}}
\newcommand{\AlertTok}[1]{\textcolor[rgb]{0.94,0.16,0.16}{#1}}
\newcommand{\AnnotationTok}[1]{\textcolor[rgb]{0.56,0.35,0.01}{\textbf{\textit{#1}}}}
\newcommand{\AttributeTok}[1]{\textcolor[rgb]{0.13,0.29,0.53}{#1}}
\newcommand{\BaseNTok}[1]{\textcolor[rgb]{0.00,0.00,0.81}{#1}}
\newcommand{\BuiltInTok}[1]{#1}
\newcommand{\CharTok}[1]{\textcolor[rgb]{0.31,0.60,0.02}{#1}}
\newcommand{\CommentTok}[1]{\textcolor[rgb]{0.56,0.35,0.01}{\textit{#1}}}
\newcommand{\CommentVarTok}[1]{\textcolor[rgb]{0.56,0.35,0.01}{\textbf{\textit{#1}}}}
\newcommand{\ConstantTok}[1]{\textcolor[rgb]{0.56,0.35,0.01}{#1}}
\newcommand{\ControlFlowTok}[1]{\textcolor[rgb]{0.13,0.29,0.53}{\textbf{#1}}}
\newcommand{\DataTypeTok}[1]{\textcolor[rgb]{0.13,0.29,0.53}{#1}}
\newcommand{\DecValTok}[1]{\textcolor[rgb]{0.00,0.00,0.81}{#1}}
\newcommand{\DocumentationTok}[1]{\textcolor[rgb]{0.56,0.35,0.01}{\textbf{\textit{#1}}}}
\newcommand{\ErrorTok}[1]{\textcolor[rgb]{0.64,0.00,0.00}{\textbf{#1}}}
\newcommand{\ExtensionTok}[1]{#1}
\newcommand{\FloatTok}[1]{\textcolor[rgb]{0.00,0.00,0.81}{#1}}
\newcommand{\FunctionTok}[1]{\textcolor[rgb]{0.13,0.29,0.53}{\textbf{#1}}}
\newcommand{\ImportTok}[1]{#1}
\newcommand{\InformationTok}[1]{\textcolor[rgb]{0.56,0.35,0.01}{\textbf{\textit{#1}}}}
\newcommand{\KeywordTok}[1]{\textcolor[rgb]{0.13,0.29,0.53}{\textbf{#1}}}
\newcommand{\NormalTok}[1]{#1}
\newcommand{\OperatorTok}[1]{\textcolor[rgb]{0.81,0.36,0.00}{\textbf{#1}}}
\newcommand{\OtherTok}[1]{\textcolor[rgb]{0.56,0.35,0.01}{#1}}
\newcommand{\PreprocessorTok}[1]{\textcolor[rgb]{0.56,0.35,0.01}{\textit{#1}}}
\newcommand{\RegionMarkerTok}[1]{#1}
\newcommand{\SpecialCharTok}[1]{\textcolor[rgb]{0.81,0.36,0.00}{\textbf{#1}}}
\newcommand{\SpecialStringTok}[1]{\textcolor[rgb]{0.31,0.60,0.02}{#1}}
\newcommand{\StringTok}[1]{\textcolor[rgb]{0.31,0.60,0.02}{#1}}
\newcommand{\VariableTok}[1]{\textcolor[rgb]{0.00,0.00,0.00}{#1}}
\newcommand{\VerbatimStringTok}[1]{\textcolor[rgb]{0.31,0.60,0.02}{#1}}
\newcommand{\WarningTok}[1]{\textcolor[rgb]{0.56,0.35,0.01}{\textbf{\textit{#1}}}}
\usepackage{graphicx}
\makeatletter
\newsavebox\pandoc@box
\newcommand*\pandocbounded[1]{% scales image to fit in text height/width
  \sbox\pandoc@box{#1}%
  \Gscale@div\@tempa{\textheight}{\dimexpr\ht\pandoc@box+\dp\pandoc@box\relax}%
  \Gscale@div\@tempb{\linewidth}{\wd\pandoc@box}%
  \ifdim\@tempb\p@<\@tempa\p@\let\@tempa\@tempb\fi% select the smaller of both
  \ifdim\@tempa\p@<\p@\scalebox{\@tempa}{\usebox\pandoc@box}%
  \else\usebox{\pandoc@box}%
  \fi%
}
% Set default figure placement to htbp
\def\fps@figure{htbp}
\makeatother
\setlength{\emergencystretch}{3em} % prevent overfull lines
\providecommand{\tightlist}{%
  \setlength{\itemsep}{0pt}\setlength{\parskip}{0pt}}
\setcounter{secnumdepth}{-\maxdimen} % remove section numbering
\usepackage{bookmark}
\IfFileExists{xurl.sty}{\usepackage{xurl}}{} % add URL line breaks if available
\urlstyle{same}
\hypersetup{
  pdftitle={Eval5\_UnixAvanzado},
  pdfauthor={OscarMacielC},
  hidelinks,
  pdfcreator={LaTeX via pandoc}}

\title{Eval5\_UnixAvanzado}
\author{OscarMacielC}
\date{2025-06-19}

\begin{document}
\maketitle

\pandocbounded{\includegraphics[keepaspectratio]{https://www.oscarmacielc.com/_next/image?url=\%2F_next\%2Fstatic\%2Fmedia\%2FOscarLogo.cdfce2e5.png&w=64&q=75}}
\href{https://github.com/OscarMacielC/RHerramientasProductividad}{El
github para la clase}
\pandocbounded{\includegraphics[keepaspectratio]{https://www.oscarmacielc.com/_next/image?url=\%2F_next\%2Fstatic\%2Fmedia\%2FOscarLogo.cdfce2e5.png&w=64&q=75}}

\begin{center}\rule{0.5\linewidth}{0.5pt}\end{center}

\subsection{Pregunta 1}\label{pregunta-1}

\textbf{¿Cuál enunciado describe las razones por las que recomendamos
usar git y GitHub al trabajar en proyectos de análisis de datos?}\\
\textbf{Respuesta:} Git y GitHub permiten un fácil control de versiones,
colaboración y compartición de recursos.

\begin{center}\rule{0.5\linewidth}{0.5pt}\end{center}

\subsection{Pregunta 2}\label{pregunta-2}

\textbf{Selecciona los pasos necesarios para:}

\begin{itemize}
\tightlist
\item
  Crear un directorio llamado \texttt{project-clone},\\
\item
  Clonar el contenido de un repositorio git en la siguiente URL en ese
  directorio (\texttt{https://github.com/user123/repo123.git}), y\\
\item
  Listar el contenido del repositorio clonado.
\end{itemize}

\textbf{Respuesta:}\\
\texttt{mkdir\ project-clone\ cd\ project-clone\ git\ clone\ https://github.com/user123/repo123.git\ ls}

\begin{center}\rule{0.5\linewidth}{0.5pt}\end{center}

\subsection{Pregunta 3}\label{pregunta-3}

\textbf{Has clonado un repositorio de GitHub y añadiste una nueva línea
al archivo \texttt{heights.txt} usando
\texttt{echo\ "165"\ \textgreater{}\textgreater{}\ heights.txt}. Luego
ejecutaste \texttt{git\ status}.}\\
\textbf{¿Qué mensaje se devuelve y qué significa?}\\
\textbf{Respuesta:}
\texttt{modified:\ heights.txt,\ no\ changes\ added\ to\ commit} → El
archivo fue modificado, pero los cambios no han sido preparados ni
confirmados en el repositorio local.

\begin{center}\rule{0.5\linewidth}{0.5pt}\end{center}

\subsection{Pregunta 4}\label{pregunta-4}

\textbf{Intentaste subir cambios usando:}

\begin{Shaded}
\begin{Highlighting}[]
\FunctionTok{git}\NormalTok{ add modified\_file.txt}
\FunctionTok{git}\NormalTok{ commit }\AttributeTok{{-}m} \StringTok{"minor changes to file"}\NormalTok{ modified\_file.txt}
\FunctionTok{git}\NormalTok{ pull}
\end{Highlighting}
\end{Shaded}

\textbf{¿Por qué no se añadió el archivo al repositorio remoto?}\\
\textbf{Respuesta:} Se debería usar \texttt{git\ push} en lugar de
\texttt{git\ pull}.

\begin{center}\rule{0.5\linewidth}{0.5pt}\end{center}

\subsection{Pregunta 5}\label{pregunta-5}

\textbf{¿Qué secuencia convierte un directorio local en un repositorio
Git y sube un archivo README?}\\
\textbf{Respuesta:}

\begin{Shaded}
\begin{Highlighting}[]
\BuiltInTok{echo} \StringTok{"A new repository with my scripts and data"} \OperatorTok{\textgreater{}}\NormalTok{ README.txt}
\FunctionTok{git}\NormalTok{ init}
\FunctionTok{git}\NormalTok{ add README.txt}
\FunctionTok{git}\NormalTok{ commit }\AttributeTok{{-}m} \StringTok{"First commit. Adding README file."}
\FunctionTok{git}\NormalTok{ remote add origin https://github.com/user123/repo123.git}
\FunctionTok{git}\NormalTok{ push}
\end{Highlighting}
\end{Shaded}

\begin{center}\rule{0.5\linewidth}{0.5pt}\end{center}

\subsection{Pregunta 6}\label{pregunta-6}

\textbf{Realizaste un cambio local, hiciste \texttt{git\ push} y
obtuviste el mensaje: \texttt{Everything\ up-to-date}.}\\
\textbf{¿Qué comando olvidaste?}\\
\textbf{Respuesta:} \texttt{git\ commit}

\begin{center}\rule{0.5\linewidth}{0.5pt}\end{center}

\subsection{Pregunta 7}\label{pregunta-7}

\textbf{Clonaste un repositorio. \texttt{git\ status} indica que todo
está actualizado, pero sabes que hubo cambios remotos.}\\
\textbf{¿Cómo sincronizas estos cambios en un solo comando?}\\
\textbf{Respuesta:} \texttt{git\ pull}

\begin{center}\rule{0.5\linewidth}{0.5pt}\end{center}

\subsection{Pregunta 8}\label{pregunta-8}

\textbf{¿Qué comando permite ocultar un archivo en Unix/Linux?}\\
\textbf{Respuesta:} Colocar un punto \texttt{.} al inicio del nombre del
archivo.

\begin{center}\rule{0.5\linewidth}{0.5pt}\end{center}

\subsection{Pregunta 9}\label{pregunta-9}

\textbf{¿Qué comando eliminará todos los archivos que comienzan con
D?}\\
\textbf{Respuesta:} \texttt{rm\ D*}

\begin{center}\rule{0.5\linewidth}{0.5pt}\end{center}

\subsection{Pregunta 10}\label{pregunta-10}

\textbf{¿Qué alias permite mostrar los 10 archivos más recientes
ordenados por fecha, mostrando nombre, tamaño y tipo?}\\
\textbf{Respuesta:}
\texttt{alias\ seetop=\textquotesingle{}ls\ -lt\ \textbar{}\ head\textquotesingle{}}

\end{document}
